\chapter{Realisierungsbewusste Optimierung}
\index{Optimierung!realisierungsbewusst}

\section*{Motivation}
Klassische Alignment-Optimierung wirkt unmittelbar auf analytische
Zielgeometrie. Reale Gleise unterscheiden sich jedoch von ihren analytischen
Zielzuständen. Das physisch realisierte Alignment wird durch Bauverfahren,
Maschinenverhalten, Schotter- und Gleismechanik, lokale Korrekturgrenzen,
Diskretisierung und Messrückkopplung beeinflusst.

Optimierung soll deshalb
\[
J\bigl(\mathcal R(\kappa,\phi)\bigr)
\]
statt \(J(\kappa,\phi)\) minimieren. In AIM optimiert sie das erwartete
realisierte operative Verhalten und nicht rein analytische Geometrie.

\section{Realisierungsoperator}
\index{Realisierung!Operator!in Optimierung}
\begin{definition}[Realisierungsoperator]
\[
\mathcal R:\mathrm{pose3}_{\mathrm{target}}
\mapsto\mathrm{pose3}_{\mathrm{real}}.
\]
\end{definition}
Der Operator stellt die Transformation vom analytischen Zielzustand zum
physisch realisierten Eisenbahnzustand dar. Er kann Abtastung,
Maschinenwirkung, Glättung, lokale Korrekturkerne und strukturelle
Gleisreaktion umfassen.

\section{Optimierung auf realisierten Zuständen}
Zielfunktionen sollen auf realisierten statt analytischen Zielzuständen
bewertet werden.
\begin{definition}[Realisierungsbewusste Optimierung]
\index{Optimierung!auf realisierten Zuständen}
\[
\min_{\kappa,\phi}\;J\bigl(\mathcal R(\kappa,\phi)\bigr).
\]
\end{definition}
Diese Formulierung integriert Realisierungseffekte ausdrücklich.

\section{pose3-Optimierung}
\index{Zustand!pose3!Optimierung}
Optimierung wirkt auf den gekoppelten pose3-Zustand
\[
(\gamma,\theta,\phi).
\]
Krümmung und Überhöhung müssen deshalb gemeinsam optimiert werden. Dies
spiegelt die dynamische Kopplung von Geometrie und Fahrzeugverhalten.

\section{Dynamische Ziele}
\subsection{Effektive Querbeschleunigung}
\begin{definition}[Effektive Querbeschleunigung]
\index{Beschleunigung!effektive Querbeschleunigung!Optimierungsziel}
\[
a_{\mathrm{eff}}=v^2\kappa-g\sin\phi.
\]
\end{definition}
Diese Größe ist die dynamisch maßgebliche Beschleunigung am Fahrzeug.

\subsection{Ruck}
\begin{definition}[Ruck]
\index{Ruck!Optimierungsziel}
\[
j=v^3\kappa'-gv\cos\phi\,\phi'.
\]
\end{definition}
Ruckbedingungen beeinflussen Übergangsentwurf und Fahrgastkomfort stark.

\subsection{Typische Zielfunktionskomponenten}
Typische Ziele sind Krümmungs- und Überhöhungsglattheit, Ruckminimierung,
Regelung des Überhöhungsfehlbetrags, Fahrzeugkomfort, Energieeffizienz und
Realisierungsrobustheit.

\section{Realisierungsbedingungen}
\index{Bedingung!Realisierung}
Optimierung muss Maschinen- und Stopfgrenzen, Schotterverhalten, lokale
Korrekturbereiche und Messauflösung einhalten. Der zulässige Entwurfsraum wird
nicht nur von Geometrie und Normen, sondern auch von Realisierbarkeit begrenzt.

\section{Inverse Realisierung}
Das Problem kann als inverses Realisierungsproblem verstanden werden.
\begin{definition}[Formulierung inverser Realisierung]
\index{Realisierung!inverse Formulierung}
\[
\mathcal R(\kappa_{\mathrm{target}},\phi_{\mathrm{target}})
\approx(\kappa_{\mathrm{desired}},\phi_{\mathrm{desired}}).
\]
\end{definition}
Der analytische Zielzustand ist damit ein Urbild unter dem
Realisierungsoperator.

\section{Frequenzinterpretation}
\index{Frequenzinterpretation!Realisierungsoperator als Tiefpassfilter}
Der Realisierungsoperator wirkt näherungsweise als Tiefpass auf Krümmung und
Überhöhung. Hochfrequente Komponenten werden gedämpft. Optimierung muss daher
geometrisch bedeutungslose hochfrequente Korrekturen vermeiden, die die
Realisierung nicht überstehen.

\section{Dünn besetzte und hybride Optimierung}
\index{Optimierung!dünn besetzt}
\index{Optimierung!hybrid}
Dünn besetzte und hybride Modelle eignen sich besonders für
realisierungsbewusste Optimierung. Dünn besetzte Strukturen bieten numerische
Stabilität, geringe Parameterdimension und interpretierbare Geometrie.
Hybride Modelle erlauben zusätzlich lokale Korrekturen,
Realisierungskompensation und adaptive Verfeinerung.

\section{Laufzeitoptimierung}
\index{Optimierung!Laufzeit}
Optimierung kann während Bau, Instandhaltung, Messrückkopplung und operativer
Überwachung kontinuierlich wirken. Sie wird damit zum Laufzeitprozess statt zu
einem rein offline ausgeführten Planungsschritt.

\section{Fahrzeugbezogene Interpretation}
\index{Fahrzeugraum!Optimierungsinterpretation}
Der physisch maßgebliche Gegenstand ist nicht zwingend die Gleismittellinie,
sondern die Bewegung des Fahrzeugs im Raum. Dazu gehören Fahrzeugschwerpunkt,
Drehgestellbewegung, Fahrzeughüllraum, Bahnsteiginteraktion und
Lichtraumverhalten. Diese Sicht verbindet realisierungsbewusste Optimierung mit
schwerpunktorientierten Alignment-Konzepten \cite{Hasslinger2002Schwerpunkt}.

\section{Ellipsoidbewusste Optimierung}
\index{Bezugssystem!ellipsoidbewusste Optimierung}
Für großräumige oder hochgenaue Systeme kann unmittelbar auf der
Erdreferenzfläche optimiert werden. Dann wird Krümmung zur geodätischen
Krümmung, Rahmen werden oberflächengebunden und World--Track-Projektion hängt
von Ellipsoidgeometrie ab.

\section{Zentrale Erkenntnis}
\begin{proposition}
Das maßgebliche Optimierungsziel ist nicht das analytische Alignment selbst,
sondern das erwartete realisierte operative Verhalten.
\end{proposition}
Dies verändert die Interpretation von Alignment-Optimierung grundlegend.

\section{Verbindung zu AIM}
\begin{summary}
Realisierungsbewusste Optimierung erweitert klassische Alignment-Optimierung
zu physisch realisierbaren Eisenbahnzuständen. In AIM wird sie auf
pose3-Zuständen, Realisierungsoperatoren, Projektionsstrukturen zur Laufzeit und
dynamisch relevantem fahrzeugbezogenem Verhalten formuliert.

Damit entsteht ein einheitlicher Rahmen für Geometrie, Realisierung, Dynamik,
Optimierung und operative Interpretation.
\end{summary}
